\section{Tasks and evaluation protocol}
\label{sec:tasks}

Two tasks share the cube, the scenario generator and one success notion.
This section defines the tasks, how scenarios are generated and seeded,
and the gating rule that keeps failed demonstrations out of every
dataset.

\paragraph{Single-arm pick-and-place.}
The cube spawns at a seeded position on one side of the table and must
be placed on a seeded target on the same side. The acting arm is the one
whose side the cube is on; sides alternate across scenarios so the
dataset covers both arms. The idle arm holds its neutral pose with the
gripper open.

\paragraph{Bimanual relay.}
The cube spawns on the left; the target lies on the right, outside the
left arm's reach. The left arm picks the cube and lays it down at a
midline staging point both arms reach, then retreats to its start; the
right arm picks the staged cube and places it on the target. The task is
genuinely bimanual --- neither arm can complete it alone --- and
sequential, which the language instruction reflects.

This design replaced a mid-air hand-over in which one arm held the
object while the other grasped it. The mid-air version failed for two
reasons we could not tune away: a single parallel-jaw grasp leaves the
object's yaw about the vertical unconstrained, so the carried object
drifted out of the second gripper's approach plane; and two gripper
bodies converging on one small object clashed geometrically. Staging on
the table restores a fully constrained object pose before the second
grasp, decomposing the relay into two instances of the already-reliable
single-arm grasp (Section~\ref{sec:difficulties}).

\paragraph{Scenario generation.}
A scenario fixes the cube spawn, the target, and (for the relay) the
staging point, all sampled from seeded uniform ranges over the reachable
table region. An analytic reachability check rejects samples either arm
cannot serve; rejected samples are resampled deterministically, so a
seed always denotes the same scenario.

\paragraph{Seed protocol.}
Three disjoint seed pools cover training (one thousand seeds, one
demonstration per seed), validation (ten seeds) and evaluation (thirty
seeds). Training, checkpoint selection and final evaluation therefore
never share a scenario. Alongside this \emph{full} mode, a \emph{simple}
mode uses a single fixed scenario for every phase: all demonstrations,
validation rollouts and evaluation trials run on the same scene. Simple
mode is the overfit-one-scenario sanity check --- a policy that cannot
master one fixed scenario has a pipeline bug, not a data problem --- and
it runs first, cheaply, before the full campaign.

\paragraph{Success and gating.}
An episode succeeds when the cube rests within \SI{2}{\centi\metre} of
the target centre, has settled (near its resting height, moving slower
than a small velocity threshold), and both grippers have released.
During collection this criterion gates every episode: a failed episode
is discarded --- never written to the dataset --- and its seed is
retried a bounded number of times before being skipped and recorded as
skipped. Before any training run, an oracle gate dry-runs both
demonstrators and aborts the experiment if the teleoperation-pipeline
oracle falls below task-specific thresholds
(Section~\ref{sec:experiments}); training on a degraded oracle would
silently poison every downstream number.
